\section*{Coverage and qualifications at a glance}
\addcontentsline{toc}{section}{Coverage and qualifications at a glance}

The starting inventory is the mathematical content of the specified
Wikipedia article, revision \texttt{1375228092}~\cite{wikipedia}.
The proofs below are developed afresh as finite algebra and finite
arithmetic. Classical background and related formulations can be found
in the Stacks Project and Conrad's notes~\cite{stacks,conrad}; their
abstract existence arguments are not imported as iPRA lemmas.

In the table, \textbf{U} means a universal finite ring identity, interpreted
arithmetically as an integer-polynomial identity; \textbf{C} means a
finite construction from the stated certificates in a fixed effective
presentation; and \textbf{Q} means a quantitative rational-approximation
version with a supplied or explicitly proved modulus. Every row refers
to a human proof in this manuscript. None denotes a completed Lean
kernel proof.

\small
\begin{longtable}{@{}p{0.07\textwidth}p{0.39\textwidth}p{0.45\textwidth}@{}}
\toprule
ID & Result family & Exact scope \\
\midrule
\endfirsthead
\toprule ID & Result family & Exact scope \\ \midrule
\endhead
\bottomrule
\endfoot
A1 & Determinant toolkit & U: signs, alternation, cofactors, adjugates, multiplicativity. \\
A2 & Characteristic coefficients & U: leading and constant coefficients, principal minors, trace coefficient. \\
A3 & Cayley--Hamilton & U: all dimensions, arbitrary commutative rings, including the zero ring. \\
A4 & Polynomial adjugate & U: every adjugate coefficient is an explicit polynomial in the matrix. \\
A5 & Evaluation and power reduction & U/C: restricted noncommutative evaluation; monic remainder and explicit power recurrence. \\
A6 & Inversion and negative powers & C: a supplied inverse of the determinant; exact numerical examples. \\
A7 & Minimal polynomials & C: bounded elimination over a fixed effective field; divisibility of the characteristic polynomial. \\
A8 & Companion matrices & U: prescribed characteristic polynomial and exact annihilator ideal. \\
A9 & Eigenvectors and roots & C: equivalence over an effective field; ring consequences with a unit component. \\
A10 & Similarity and split spectra & U/C: supplied inverses, triangularization or factorization; coefficient and power-trace formulas. \\
T1 & Jacobi and adjugate traces & U: finite formal differentiation. \\
T2 & Newton identities & U: all degrees, without division or eigenvalues. \\
T3 & Bell-polynomial formulas & U: factorial-multiplied identities; division only with justified units or exact integer quotients. \\
T4 & Hessenberg determinant formula & U: determinant equals $k!e_k$; explicit formulas through order four. \\
T5 & Faddeev--LeVerrier & C: correctly initialized; supplied units or proved exact integer division. \\
T6 & Exponential/logarithm and partitions & C: every finite formal precision, with integer denominators invertible. \\
M1 & Finitely generated modules & C: supplied generator images and coordinates; no basis required. \\
M2 & Ideal coefficients and Nakayama & C: finite membership and unit certificates. \\
M3 & Closure of integral elements & C: monic input equations give a degree-$\prod_jd_j$ output equation. \\
M4 & Algebraic-number minimal polynomials & C: finite multiplication table; integral coefficients for integral elements. \\
I1 & Lagrange and Sylvester formulas & C: supplied roots with unit differences; spectral idempotents. \\
I2 & Hermite interpolation & C: repeated-root jets, Hasse derivatives, explicit truncated inverses. \\
F1 & Polynomial matrix-function reduction & Q: primitive recursive rational coefficients and weighted finite-tail bound. \\
F2 & Exponential, sine, cosine & Q: rational matrix and parameter inputs, explicit factorial cutoff. \\
F3 & Rotation formulas & Q: planar, Pauli, and Rodrigues identities from supplied polynomial relations. \\
Q1 & Quaternion and split-quaternion matrices & C: faithful representations give degree $2n$ annihilators; also degree $4n$ regular representation. \\
Q2 & One octonion & U/C: quadratic norm identity, with rational coordinate interpretation. \\
\end{longtable}
\normalsize

The unrestricted analytic functional calculus is not claimed as a literal
first-order iPRA theorem. Neither are a root-finding or Jordan-form
existence theorem for arbitrary fields, density of diagonalizable complex
matrices, or the modified matrix-octonion theorem. The article's broader
Lie-group references and its exterior-algebra and combinatorial proof
methods are not developed here. The main theorem already has two
complete finite proofs; those other methods are unnecessary dependencies.

Several source formulations require correction. Over a ring, an inverse
requires a unit determinant. Trace expressions with factorial
denominators must be qualified. The stated partition product has an
extra sign, and the displayed Faddeev--LeVerrier recursion lacks its
first valid initialization. A nonpolynomial analytic function cannot be
written $qp+r$ with both $q,r$ polynomials. These points are proved or
diagnosed in the appropriate sections rather than used as premises.

